\documentclass[12pt]{article}
\usepackage[a4paper,margin=1in]{geometry}
\usepackage[utf8]{inputenc}
\usepackage[T1]{fontenc}
\usepackage{amsmath,amssymb}
\usepackage{textcomp}
\usepackage{parskip}
\begin{document}
\begin{center}
{\Large \textbf{2013 STEP I}}
\end{center}

\subsection*{Page 1}
\begin{flushleft}
\bigskip\textbf{Section A: Pure Mathematics}\par
\bigskip\textbf{Question 1.} (i) Use the substitution \sqrtx = y (where y \leq 0) to find the real root of the equation\par
\[x + 3 \sqrt\{x\} 1 = 0 .\]
- 2\par
\textbf{(ii)} Find all real roots of the following equations:\par
\[(a) x + 10 \sqrt\{x\} + 2 22 = 0 ;\]
-\par
\[(b) x^2 4x + \sqrt\{2x2\} 8x 3 9 = 0 .\]
- - - -\par
\bigskip\textbf{Question 2.} In this question, x denotes the greatest integer that is less than or equal to x, so that\par
\[\geq (\]
\[2.9 = 2 = 2.0 and 1.5 = 2.\]
\[\geq ( \geq ( \geq- ( -\]
x\par
\[The function f is defined, for x = 0, by f(x) = \geq ( .\]
) x\par
\textbf{(i)} Sketch the graph of y = f(x) for 3 \geq x \geq 3 (with x = 0).\par
- )\par
\textbf{(ii)} By considering the line y = 7 on your graph, or otherwise, solve the equation f(x) = 7 .\par
12 12\par
\[Solve also the equations f(x) = 17 and f(x) = 4 .\]
24 3\par
\textbf{(iii)} Find the largest root of the equation f(x) = 9 .\par
Give necessary and sufficient conditions, in the form of inequalities, for the equation f(x) = c\par
\[to have exactly n roots, where n \leq 1.\]
\bigskip\textbf{Question 3.} For any two points X and Y , with position vectors x and y respectively, X Y is defined to\par
∗\par
be the point with position vector λx + (1 λ)y, where λ is a fixed number.\par
-\par
\textbf{(i)} If X and Y are distinct, show that X Y and Y X are distinct unless λ takes a certain\par
∗ ∗\par
value (which you should state).\par
\textbf{(ii)} Under what conditions are (X Y ) Z and X (Y Z) distinct?\par
∗ ∗ ∗ ∗\par
\textbf{(iii)} Show that, for any points X, Y and Z,\par
\[(X Y ) Z = (X Z) (Y Z)\]
∗ ∗ ∗ ∗ ∗\par
and obtain the corresponding result for X (Y Z).\par
∗ ∗\par
\textbf{(iv)} The points P 1, P 2, . . . are defined by P 1 = X Y and, for n \leq 2, P n = P n-1 Y . Given\par
∗ ∗\par
\[that X and Y are distinct and that 0 < λ < 1, find the ratio in which P divides the\]
n\par
line segment XY .\par
\end{flushleft}
\subsection*{Page 2}
\begin{flushleft}
\bigskip\textbf{Question 4.} (i) Show that, for n > 0,\par
\[1 4 \pi tann x sec 2 x dx = 1 and 4 1 \pi secnx tan x dx = (\sqrt\{2\})n - 1 .\]
\[\leq0 n + 1 \leq0 n\]
\textbf{(ii)} Evaluate the following integrals:\par
\[1 \pi 1 \pi\]
4 4 4 2 2\par
x sec x tan x dx and x sec x tan x dx .\par
\[\leq0 \leq0\]
\bigskip\textbf{Question 5.} The point P has coordinates (x, y) which satisfy\par
2 2\par
\[x + y + kxy + 3x + y = 0 .\]
\textbf{(i)} Sketch the locus of P in the case k = 0, giving the points of intersection with the\par
coordinate axes.\par
\textbf{(ii)} By factorising 3x2 + 3y2 + 10xy, or otherwise, sketch the locus of P in the case k = 10 ,\par
giving the points of intersection with the coordinate axes.\par
\textbf{(iii)} In the case k = 2, let Q be the point obtained by rotating P clockwise about the origin\par
by an angle θ, so that the coordinates (X, Y ) of Q are given by\par
\[X = x cos θ + y sin θ , Y = x sin θ + y cos θ .\]
-\par
\[Show that, for θ = 45 ^\circ , the locus of Q is \sqrt\{2Y\} = (\sqrt\{2X\} + 1)2 1.\]
-\par
\[Hence, or otherwise, sketch the locus of P in the case k = 2, giving the equation of the\]
line of symmetry.\par
\bigskip\textbf{Question 6.} By considering the coefficient of xr in the series for (1 + x)(1 + x)n, or otherwise, obtain the\par
following relation between binomial coefficients:\par
n n n + 1\par
\[+ = (1 \geq r \geq n).\]
\[\geqr( \geqr 1( \geq r (\]
-\par
The sequence of numbers B 0, B 1, B 2, . . . is defined by\par
m m\par
2m j 2m + 1 k\par
\[B 2m = - and B 2m+1 = - .\]
\[\geq j ( \geq k (\]
\[)j=0 )k=0\]
\[Show that B n+2 B n+1 = B n (n = 0, 1, 2, . . . ).\]
-\par
What is the relation between the sequence B 0, B 1, B 2, . . . and the Fibonacci sequence F 0, F 1,\par
\[F 2, . . . defined by F 0 = 0, F 1 = 1 and F n = F n-1 + F n-2 for n \leq 2?\]
\end{flushleft}
\subsection*{Page 3}
\begin{flushleft}
\bigskip\textbf{Question 7.} (i) Use the substitution y = ux, where u is a function of x, to show that the solution of the\par
differential equation\par
dy x y\par
\[= + (x > 0, y > 0)\]
dx y x\par
\[that satisfies y = 2 when x = 1 is\]
\[y = x \sqrt\{4\} + 2 ln x (x > e -2 ).\]
\textbf{(ii)} Use a substitution to find the solution of the differential equation\par
dy x 2y\par
\[= + (x > 0, y > 0)\]
dx y x\par
\[that satisfies y = 2 when x = 1.\]
\textbf{(iii)} Find the solution of the differential equation\par
\[dy x^2 2y\]
\[= + (x > 0, y > 0)\]
dx y x\par
\[that satisfies y = 2 when x = 1.\]
\bigskip\textbf{Question 8.} (i) The functions a, b, c and d are defined by\par
\[a(x) = x^2 ( < x < ),\]
-\infty \infty\par
\[b(x) = ln x (x > 0),\]
\[c(x) = 2x ( < x < ),\]
-\infty \infty\par
\[d(x) = \sqrt\{x\} (x \leq 0) .\]
Write down the following composite functions, giving the domain and range of each:\par
cb, ab, da, ad.\par
\textbf{(ii)} The functions f and g are defined by\par
\[f(x) = \sqrt\{x^2\} 1 ( x \leq 1),\]
- | |\par
\[g(x) = \sqrt\{x^2\} + 1 ( < x < ).\]
-\infty \infty\par
Determine the composite functions fg and gf, giving the domain and range of each.\par
\textbf{(iii)} Sketch the graphs of the functions h and k defined by\par
\[h(x) = x + \sqrt\{x^2\} 1 (x \leq 1),\]
-\par
\[k(x) = x \sqrt\{x^2\} 1 ( x \leq 1),\]
- - | |\par
justifying the main features of the graphs, and giving the equations of any asymptotes.\par
Determine the domain and range of the composite function kh.\par
\end{flushleft}
\subsection*{Page 4}
\begin{flushleft}
\bigskip\textbf{Section B: Mechanics}\par
\bigskip\textbf{Question 9.} Two particles, A and B, are projected simultaneously towards each other from two points\par
which are a distance d apart in a horizontal plane. Particle A has mass m and is projected at\par
speed u at angle α above the horizontal. Particle B has mass M and is projected at speed v\par
at angle β above the horizontal. The trajectories of the two particles lie in the same vertical\par
plane.\par
The particles collide directly when each is at its point of greatest height above the plane.\par
Given that both A and B return to their starting points, and that momentum is conserved in\par
the collision, show that\par
\[m cot α = M cot β .\]
Show further that the collision occurs at a point which is a horizontal distance b from the\par
point of projection of A where\par
M d\par
\[b = ,\]
m + M\par
and find, in terms of b and α, the height above the horizontal plane at which the collision\par
occurs.\par
\bigskip\textbf{Question 10.} Two parallel vertical barriers are fixed a distance d apart on horizontal ice. A small ice\par
hockey puck moves on the ice backwards and forwards between the barriers, in the direction\par
perpendicular to the barriers, colliding with each in turn. The coefficient of friction between\par
the puck and the ice is μ and the coefficient of restitution between the puck and each of the\par
barriers is r.\par
The puck starts at one of the barriers, moving with speed v towards the other barrier. Show\par
that\par
2 2 2 2\par
\[v r v = 2r μgd\]
i+1 i\par
- -\par
where v is the speed of the puck just after its ith collision.\par
i\par
The puck comes to rest against one of the barriers after traversing the gap between them n\par
\[v^2\]
\[times. In the case r = 1, express n in terms of r and k, where k = . If r = e\]
-1\par
(where\par
) 2μgd\par
e is the base of natural logarithms) show that\par
\[n = 1 ln 1 + k(e 2 1) .\]
2 -\par
\[\int \sum\]
\[Give an expression for n in the case r = 1.\]
\end{flushleft}
\subsection*{Page 5}
\begin{flushleft}
A\par
B\par
\[1 2 \pi - α C α\]
The diagram shows a small block C of weight W initially at rest on a rough horizontal surface.\par
The coefficient of friction between the block and the surface is μ. Two light strings, AC and\par
BC, are attached to the block, making angles 1 \pi α and α to the horizontal, respectively.\par
The tensions in AC and BC are T sin β and T c 2 os β - respectively, where 0 < α + β < 1 \pi.\par
\textbf{(i)} In the case W > T sin(α + β), show that the block will remain at rest provided\par
\[W sin λ \leq T cos(α + β λ) ,\]
-\par
\[where λ is the acute angle such that tan λ = μ.\]
\textbf{(ii)} In the case W = T tan φ, where 2φ = α + β, show that the block will start to move in a\par
direction that makes an angle φ with the horizontal.\par
\end{flushleft}
\subsection*{Page 6}
\begin{flushleft}
\bigskip\textbf{Section C: Probability and Statistics}\par
12 Each day, I have to take k different types of medicine, one tablet of each. The tablets are\par
identical in appearance. When I go on holiday for n days, I put n tablets of each type in a\par
container and on each day of the holiday I select k tablets at random from the container.\par
\textbf{(i)} In the case k = 3, show that the probability that I will select one tablet of each type on\par
the first day of a three-day holiday is .\par
Write down the probability that I will be left with one tablet of each type on the last\par
day (irrespective of the tablets I select on the first day).\par
\textbf{(ii)} In the case k = 3, find the probability that I will select one tablet of each type on the\par
first day of an n-day holiday.\par
\textbf{(iii)} In the case k = 2, find the probability that I will select one tablet of each type on each\par
day of an n-day holiday, and use Stirling's approximation\par
n n\par
\[n! \sqrt\{2n\}\pi\]
\approx e\par
( )\par
to show that this probability is approximately 2\par
\[-n\sqrt\{n\}\pi\]
.\par
13 From the integers 1, 2, . . . , 52, I choose seven (distinct) integers at random, all choices being\par
equally likely. From these seven, I discard any pair that sum to 53. Let X be the random\par
variable the value of which is the number of discarded pairs. Find the probability distribution\par
\[of X and show that E(X) = 7 .\]
\[Note: 7 17 47 = 5593.\]
\[\times \times\]
\end{flushleft}
\end{document}
